\LectureAnswerSection

\begin{drillsolution}{n1d1}\textbf{先求层间法向力与共同水平加速度}\par\nobreak
图\ref{n1d1q}中上下两物体竖直加速度都为 $a_0$。设层间支持力 $N$，对上层写 $N-mg=ma_0$，故 $N=m(g+a_0)$。水平方向地面光滑，整体 $F=(M+m)a_x$；上层需要的摩擦为 $f=ma_x=mF/(M+m)$。
\fig{n1d1s}{水平加速所需摩擦与竖直运动决定的上限}{\begin{tikzpicture}
\fill(0,0)circle(2pt);\draw[force,blue](0,0)--(0,1.4)node[above]{$N=m(g+a_0)$};\draw[force,red](0,0)--(0,-1)node[below]{$mg$};\draw[force](0,0)--(1.5,0)node[right]{$f=ma_x$};\node at(3,1){$f\le\mu_sN$};
\end{tikzpicture}}


\par\medskip\Needspace{4\baselineskip}\noindent\textbf{比较所需摩擦与上限，核查失重极限}\par\nobreak
由图\ref{n1d1s}所示上限，$mF/(M+m)\le\mu_s m(g+a_0)$，得到 $F_{\max}=\mu_s(M+m)(g+a_0)$。地板支持力 $R$ 对整体提供竖直加速度，$R-(M+m)g=(M+m)a_0$，故 $R=(M+m)(g+a_0)$。若将 $a_0$ 改成向下的有符号值，接近 $-g$ 时支持和可用静摩擦均趋于零，模型在失重处需要另行判断接触。
\end{drillsolution}

\begin{drillsolution}{n1d2}\textbf{由秤读数恢复张力与两悬物方程}\par\nobreak
图\ref{n1d2q}中的理想滑轮静止，故 $T=S/2$。两悬物沿各自运动方向的方程相加给 $a=(m_1-m_2)g/M$，这里先约定 $m_1\ge m_2$。代回任一物体可得 $S=4m_1m_2g/M$。
\fig{n1d2s}{用质量差表示不对称程度}{\begin{tikzpicture}
\draw[axis](0,0)--(4.5,0)node[right]{$\delta/M$};\draw[axis](0,0)--(0,2.2)node[above]{$S/(Mg)$};\draw[world,domain=0:1,samples=40]plot({4*\x},{2*(1-\x*\x)});\node[left]at(0,2){$1$};\node[below]at(4,0){$1$};\node at(2.2,1.6){$1-(\delta/M)^2$};
\end{tikzpicture}}


\par\medskip\Needspace{4\baselineskip}\noindent\textbf{用总质量和质量差解出两质量}\par\nobreak
令质量差 $\delta=m_1-m_2\ge0$。由 $M=m_1+m_2$，先计算
\[
 M^2-\delta^2=(m_1+m_2)^2-(m_1-m_2)^2=4m_1m_2.
\]
按图\ref{n1d2s}将它代入读数式 $S=4m_1m_2g/M$，得 $SM/g=M^2-\delta^2$，即 $\delta^2=M^2-SM/g$。质量差取非负根，得到 $\delta=M\sqrt{1-S/(Mg)}$。再联立和、差，$m_1=(M+\delta)/2$、$m_2=(M-\delta)/2$，而加速度大小为 $a=g\delta/M$，所以
\[
 m_{1,2}=\frac M2\left(1\pm\sqrt{1-\frac S{Mg}}\right),\qquad
 a=g\sqrt{1-\frac S{Mg}}.
\]


\par\medskip\Needspace{4\baselineskip}\noindent\textbf{检验读数的允许范围}\par\nobreak
两质量均为正时须 $0<S\le Mg$，等号对应两质量相等；$S=0$ 仅为其中一质量趋于零的极限。测得 $S>Mg$ 则与该理想模型不相容，应检查滑轮质量、秤的运动或其他外力。
\end{drillsolution}

\begin{drillsolution}{n1d3}\textbf{按圆周曲率列径向方程求张力}\par\nobreak
图\ref{n1d3q}的速度沿切线，方向随时间改变。由圆周几何，径向加速度大小 $a_r=v^2/l$，指向悬点。向悬点列式为 $T-mg\cos\theta=mv^2/l$，从而 $T=m(g\cos\theta+v^2/l)$。
\fig{n1d3s}{加速度分为指向悬点和指向最低点的两个分量}{\begin{tikzpicture}
\fill(0,0)circle(2pt);\draw[flow,blue](0,0)--(-.65,1.15)node[above]{$v^2/l$};\draw[flow,red](0,0)--(-1.3,-.75)node[left]{$g\sin\theta$};\draw[guide](-1,-.58)--(.5,.29);\node at(2,1){径向不是零};
\end{tikzpicture}}


\par\medskip\Needspace{4\baselineskip}\noindent\textbf{沿切线求加速度，再合成完整矢量}\par\nobreak
图\ref{n1d3s}的切向不含张力，重力沿向最低点方向的分量为$mg\sin\theta$，所以$a_t=g\sin\theta$。完整加速度为这两个正交分量的矢量和，大小$\sqrt{v^4/l^2+g^2\sin^2\theta}$。只有$v=0$时径向分量才消失；把静止释放的加速度套到摆动中间时刻，会遗漏曲率导致的方向变化。还须检查$T=m(g\cos\theta+v^2/l)\ge0$。若给定角度、速率使它为负，轻绳不能提供所需的压力，绳已松弛，应改按仅受重力的自由运动处理，不能继续使用半径为$l$的圆周约束。
\end{drillsolution}

\begin{drillsolution}{n1d4}\textbf{由斜面平移确定接触所需的法向加速度}\par\nobreak
图\ref{n1d4q}沿斜面向右上取单位矢量 $\bm t=(\cos\alpha,\sin\alpha)$，向左上法向为 $\bm n=(-\sin\alpha,\cos\alpha)$。斜面加速度为 $-a_0\bm e_x$，法向投影为 $a_0\sin\alpha$。直面接触要求物块具有相同法向加速度，故 $N-mg\cos\alpha=ma_0\sin\alpha$。
\fig{n1d4s}{相对加速度只沿直斜面，法向由斜面平移确定}{\begin{tikzpicture}
\draw[axis](0,0)--(1.7,1)node[right]{$\bm t$};\draw[axis](0,0)--(-.9,1.5)node[above]{$\bm n$};\draw[flow](0,0)--(-1.8,0)node[left]{$-a_0\bm e_x$};\draw[force,red](0,0)--(0,-1.3)node[below]{$mg$};
\end{tikzpicture}}


\par\medskip\Needspace{4\baselineskip}\noindent\textbf{求支持力并相减得到沿面相对加速度}\par\nobreak
因此 $N=m(g\cos\alpha+a_0\sin\alpha)>0$。沿 $\bm t$ 的对地加速度由重力给出 $\bm a_m\cdot\bm t=-g\sin\alpha$。根据图\ref{n1d4s}减去斜面沿面分量 $-a_0\cos\alpha$，得相对沿面加速度 $a_{\rm rel}=a_0\cos\alpha-g\sin\alpha$。

当 $a_0=g\tan\alpha$ 且初始相对速度也为零时，物块才能始终相对静止；只令相对加速度为零而不给初始速度条件，仍可能匀速沿面滑动。
\end{drillsolution}
